(1) 相同费米子的总波函数在交换粒子1、2时必须变号。两粒子在最低轨道 \(\phi_0\) 的空间波函数是对称的 \(\Phi_{00}(x_1,x_2)=\phi_0(x_1)\phi_0(x_2)\)，故自旋只能取反对称单重态
\[
|s\rangle=\frac1{\sqrt2}(|\uparrow\downarrow\rangle-|\downarrow\uparrow\rangle).
\]
基态总能量为 \(E_{\rm ground}=2\varepsilon_0=\hbar\omega\)。它是唯一的基态（忽略外部规范简并），不能让两颗相同自旋的费米子都占据同一空间自旋轨道。

(2) 第一激发层由一个粒子占据 \(n=0\)、另一个占据 \(n=1\) 组成，总能量 \(E_{01}=\varepsilon_0+\varepsilon_1=2\hbar\omega\)。两个归一化空间组合是
\[
\Phi_S=\frac{\phi_0(x_1)\phi_1(x_2)+\phi_1(x_1)\phi_0(x_2)}{\sqrt2},
\quad
\Phi_A=\frac{\phi_0(x_1)\phi_1(x_2)-\phi_1(x_1)\phi_0(x_2)}{\sqrt2}.
\]
\(\Phi_S\) 对称，须配反对称自旋单重态；\(\Phi_A\) 反对称，须配三个对称自旋三重态
\[
|\uparrow\uparrow\rangle,\quad
(|\uparrow\downarrow\rangle+|\downarrow\uparrow\rangle)/\sqrt2,\quad
|\downarrow\downarrow\rangle.
\]
故该能量层有 \(1+3=4\) 个独立总态，均满足总交换反对称性。

(3) 对三重态，其空间因子为 \(\Phi_A\)，在 \(x_1=x_2=x\) 时
\[
\Phi_A(x,x)=\frac{\phi_0(x)\phi_1(x)-\phi_1(x)\phi_0(x)}{\sqrt2}=0.
\]
因此同位联合概率密度严格为零，表现为交换反相关。单重态配的是 \(\Phi_S\)，其同位振幅为 \(\Phi_S(x,x)=\sqrt2\,\phi_0(x)\phi_1(x)\)，除轨道节点外可非零。二者的区别来自空间宇称交换对称性，而非不同的直接作用势能。

(4) 如果两个粒子的自旋均固定在同一极化方向，总自旋因子对称，空间波函数必须反对称。最低允许占据是一个 \(n=0\)、另一个 \(n=1\)，故最低总能量为
\[
E_{\rm polarised}=\varepsilon_0+\varepsilon_1=2\hbar\omega.
\]
相比允许自旋单重态的两粒子基态，其能量更高，但哈密顿量中从未加入经典粒子间排斥项；这正是 Pauli 对称性约束导致的纯交换效应。
